\DocumentMetadata{
  pdfversion=2.0,
  pdfstandard=ua-2,
  lang=en-US,
  tagging=on,
  tagging-setup={math/setup=mathml-SE,tabsorder=structure}
}
\documentclass[11pt,letterpaper]{article}
\usepackage[margin=0.68in,includefoot]{geometry}
\usepackage{newcomputermodern}
\usepackage{amsmath,amscd,mathtools}
\usepackage{tikz,adjustbox}
\providecommand{\square}{\mdlgwhtsquare}
\usepackage[hidelinks]{hyperref}
\hypersetup{
  pdftitle={Logic PhD Exam, January 2009},
  pdfauthor={University of Florida Department of Mathematics},
  pdfsubject={Graduate mathematics examination},
  pdfdisplaydoctitle=true,
  bookmarksopen=true
}
\setcounter{secnumdepth}{0}
\setlength{\parindent}{0pt}
\setlength{\parskip}{0.28em}
\setlength{\emergencystretch}{3em}
\setlength{\leftmargini}{2.15em}
\setlength{\leftmarginii}{2.65em}
\setlength{\leftmarginiii}{3em}
\renewcommand{\labelenumi}{\textbf{\theenumi.}}
\pagestyle{empty}
\raggedbottom
\allowdisplaybreaks
\newenvironment{examproblems}[1]{%
  \begin{enumerate}%
  \setcounter{enumi}{#1}%
  \setlength{\topsep}{0.35em}%
  \setlength{\partopsep}{0pt}%
  \setlength{\itemsep}{0.48em}%
  \setlength{\parsep}{0.18em}%
}{\end{enumerate}}
\newenvironment{examsubparts}{%
  \begin{enumerate}%
  \setlength{\topsep}{0.18em}%
  \setlength{\partopsep}{0pt}%
  \setlength{\itemsep}{0.16em}%
  \setlength{\parsep}{0.1em}%
}{\end{enumerate}}
\newenvironment{examinstructions}{%
  \par\begingroup\itshape
}{%
  \par\endgroup\vspace{0.18em}
}
\makeatletter
\renewcommand\section{\@startsection{section}{1}{0pt}%
  {-1.25ex plus -.25ex minus -.1ex}%
  {0.55ex plus .1ex}%
  {\normalfont\large\bfseries}}
\renewcommand{\@maketitle}{%
  \newpage\null\vskip -1em%
  {\footnotesize\bfseries
    \href{https://www.ufl.edu/}{University of Florida}, \href{https://math.ufl.edu/}{Department of Mathematics}\par}%
  \vskip -0.5em%
  \tagstructbegin{tag=Title}%
    {\LARGE\bfseries\href{https://gma.math.ufl.edu/exams/logic-phd/}{Logic PhD Exam}, January 2009\par}%
  \tagstructend%
  \vskip 0.25em%
  \tagmcbegin{artifact}%
    \hrule height 1pt%
  \tagmcend%
  \vskip 1em%
}
\makeatother
\title{Logic PhD Exam, January 2009}
\author{}
\date{}
\begin{document}
\maketitle
\thispagestyle{empty}
\begin{examinstructions}
Answer six questions and at least one from each section.
\end{examinstructions}
\section{Section 1}
\begin{examproblems}{0}
\item
Sketch the proof of Gödel’s Completeness Theorem for Predicate Logic and explain the use of Henkin constants (or Skolem functions).
\item
Show that for any elementary chain \(\{\mathcal A_i\}_{i\in\omega}\) of structures, \(\mathcal A_i\) is an elementary submodel of the union \(\mathcal A\).
\item
Show that the theory of dense linear orderings without end points is complete and decidable.
\end{examproblems}
\section{Section 2}
\begin{examproblems}{3}
\item
Suppose \(2\leq\kappa\leq\lambda\) and~\(\lambda\) is infinite. Show that \(2^\lambda=\kappa^\lambda\).
\item
State and prove König’s Lemma. (Hint: trees.)
\item
Show that for any notion of forcing \(\mathcal P\) and any countable set \(\mathcal D\) of \(\mathcal P\)-dense sets, there exists a \(\mathcal D\)-generic \(\mathcal P\)-filter.
\end{examproblems}
\section{Section 3}
\begin{examproblems}{6}
\item
Sketch a proof of Gödel’s incompleteness Theorem for Arithmetic.
\item
Define many-one and truth-table reducibility and show that one implies the other but the converse does not hold in general.
\item
State the Friedberg–Muchnik Theorem and sketch the proof. Explain how the “priority” argument deals with “injury”.
\end{examproblems}
\end{document}
